% Section 2 — Related software and data ecosystem. Replaces the ECCV
% related-work section with a software/community survey relevant to
% atmospheric scientists.

We survey the existing software and data landscape for European
PM\textsubscript{2.5} forecasting along three axes: physics-based
operational models, machine-learning weather and chemistry models, and
satellite-derived high-resolution analyses. Table~\ref{tab:related}
summarises licence, native resolution, lead time and accelerator support.

\subsection{Physics-based operational models}

The Copernicus Atmosphere Monitoring Service regional ensemble (CAMS,
\citealp{cams}) blends nine national chemistry transport models
(EMEP, LOTOS-EUROS, CHIMERE, MOCAGE, SILAM, MATCH, MINNI, EURAD-IM,
GEM-AQ) at a nominal resolution of 0.1$^\circ$. Operational outputs are
freely available through the Atmosphere Data Store but the source
code of individual models is not uniformly open. EMEP-MSC-W
\citep{simpson2012emep} and LOTOS-EUROS \citep{schaap2008lotos} are
two notable exceptions, with permissive open-source licences and
active community releases. Their compute cost prohibits routine
1\,km production: a single 0.1$^\circ$ EMEP run over Europe takes
on the order of $10^4$ CPU-hours per simulated year, and a 1\,km
configuration would scale roughly with the inverse square of grid
spacing.

\subsection{ML-based weather and air-quality models}

The recent generation of data-driven weather models (Pangu-Weather,
\citealp{bi2023pangu}; GraphCast, \citealp{lam2023graphcast};
Aurora, \citealp{bodnar2024aurora}; ClimaX, \citealp{nguyen2023climax})
has demonstrated that ViT- or graph-based architectures can match or
beat operational NWP at coarse-to-medium resolution and a fraction of
the inference cost. Most of these models are released as research
code with single-GPU inference; only Aurora ships pre-trained weights
for an air-quality variant, and at a 25\,km nominal resolution. None
target the 1\,km surface-PM forecasting problem.

In the dedicated AQ literature, ConvLSTM \citep{shi2015convlstm},
SimVP \citep{gao2022simvp} and Earthformer \citep{gao2022earthformer}
have been used as fine-resolution PM\textsubscript{2.5} forecasters
in regional studies. They typically tile the domain at a single
resolution, which limits either the spatial detail (when chosen
coarse) or the receptive field (when chosen fine). CRAN-PM is to our
knowledge the first model that operates simultaneously on 25\,km
meteorology and 1\,km PM\textsubscript{2.5} via a learnt
cross-resolution attention bridge.

\subsection{High-resolution satellite-derived analyses}

The GHAP daily PM\textsubscript{2.5} product \citep{wei2023ghap}
provides 1\,km global retrospective analyses by combining MAIAC
aerosol optical depth, surface stations and meteorological reanalysis
through a tree-based ML model. Latency is on the order of two weeks,
which makes GHAP a reference for retrospective evaluation but not a
forecast product. CRAN-PM uses GHAP as both training target and as
the high-resolution input channel at the forecast initialisation
time.

\subsection{A note on WRF-Chem and the choice of physics reference}

WRF-Chem \citep{grell2005wrfchem} is the standard online-coupled
chemistry-transport extension of WRF and would in principle be the
most natural physics-based reference for a paper on PM\textsubscript{2.5}
forecasting. We do not include it in the comparison because no
operational European WRF-Chem run at our test resolution is
publicly archived for the test year 2022. Re-running WRF-Chem over
the European domain at 1\,km for a representative subset of days
would require on the order of $10^5$ CPU-hours, beyond the budget of
this software paper. CAMS regional, which bundles nine
chemistry-transport models including WRF-class physics, is therefore
used as the operational physics reference throughout this manuscript.

\subsection{Position of CRAN-PM v1.0}

CRAN-PM v1.0 sits in a previously empty cell of the matrix in
Table~\ref{tab:related}: it is the only open-source, dual-backend
(CUDA + ROCm), 1\,km-native, short-range PM\textsubscript{2.5}
forecasting tool aimed at atmospheric scientists. It does not
replace operational chemistry-transport models, which remain the
authoritative source of process understanding, but complements them
by delivering high-resolution short-range forecasts at orders of
magnitude lower compute cost.

% TODO: typeset Table~\ref{tab:related} in tables/table_related.tex
% (a software-ecosystem comparison: model name, license, native
% resolution, lead time, accelerator support, training data
% availability).
